\section{Cost and performance}
\label{sec:cost}

\status{Uniform-mesh rates are \textbf{measured} at both resolutions. The
refinement gain is bounded by two measured numbers --- cell count and remeshing
overhead --- but the production comparison has not been run.}

\subsection{Uniform mesh}

\begin{table}[htbp]
\centering
\caption{Measured cost in microseconds per cell-step, single core, uniform mesh.
Same mesh and same conditions for both methods.}
\label{tab:cost}
\small
\pgfplotstabletypeset[
  col sep=space, string type,
  columns/caso/.style={column name={Rheology}, column type={l}},
  columns/vof/.style={column name={VOF}, column type={r}},
  columns/ft/.style={column name={FT}, column type={r}},
  every head row/.style={before row=\toprule, after row=\midrule},
  every last row/.style={after row=\bottomrule},
]{\dat{custo.dat}}
\end{table}

The two methods cost the same to within five percent, and which is cheaper
depends on the resolution: front-tracking is 2.5\% more expensive at
$25\times25\times50$ and 4.9\% cheaper at $50\times50\times100$ (19.5 against
20.5 hours of wall clock for the full 1501 steps).

The weak version of that statement is the defensible one, and it is still worth
making, because the strong opposite is often assumed: an explicit surface sounds
much more expensive than a scalar field, and it is not. VOF pays for the PLIC
reconstruction and the split advection in every cell near the interface;
front-tracking pays for advecting a few thousand vertices and rebuilding one
spatial index. Those turn out to be comparable.

\paragraph{How not to measure this.}
An earlier version of this table reported front-tracking as 12\% cheaper. That
figure came from eleven-step runs, where setup is a large fraction of the total
--- and, it turned out, a \emph{different} fraction for the two methods. Only
the complete runs give a usable ratio.

\paragraph{The cost per cell-step rises with problem size.}
By 14\% for VOF and 6\% for front-tracking between the two resolutions. Part of
that is cache pressure and part is that the two reference runs shared the
machine on separate cores while the coarse runs did not; these measurements do
not separate the two causes. Every rate in this paper states the problem size
and whether the machine was shared, because a rate without those is not
transferable --- an early estimate in this work extrapolated a small-problem
rate to the large problem and underpredicted the wall clock by a factor of about
1.6.

\subsection{Output cost}

Worth a paragraph because it dominated a first attempt at the reference
resolution. The field output is an unstructured grid with eight unshared corner
points per cell: $250\,000$ cells become two million points, about 300~MB per
frame, roughly 20 minutes to write. At 120 frames that is some 40 hours of
writing against 16 hours of computation. The extracted interface is a factor of
226 smaller and already carries the interpolated fields at its vertices, so the
field output is written at a coarse cadence and the interface at every frame.

\subsection{Adaptive against uniform}

The remeshing machinery works (Section~\ref{sec:amr}); what is missing is a
production run on an adapted mesh. Two measured numbers already bound what to
expect.

\paragraph{Cell count.} For the reference geometry with a minimum of five fine
cells per side, the criterion produces, against the equivalent uniform fine
mesh:

\begin{center}
\small
\begin{tabular}{llrrr}
\toprule
Base & Levels & Criterion & Uniform fine & Factor \\
\midrule
$50\times50\times100$ & 1 & 379\,136 & 2\,000\,000 & 5.3 \\
$50\times50\times100$ & 2 & 1\,128\,584 & 16\,000\,000 & 14.2 \\
$50\times50\times100$ & 3 & 5\,458\,224 & 128\,000\,000 & 23.5 \\
\bottomrule
\end{tabular}
\end{center}

Base 50 with two levels is an effective $200\times200\times400$ at 1.13 million
cells. At the measured rates that is of order five days against roughly seventy
for the equivalent uniform mesh --- which is what makes the fine resolution
reachable at all.

\paragraph{Remeshing overhead.} The criterion costs 0.008~s; the reconstruction
costs between 1.9 and 4.7~s, varying with how much the mesh changed. At one
remesh every 25 steps that is under 1\% of a step. The bottleneck is the
reconstruction, not the criterion, by a factor of 230.

\paragraph{The band over-delivers, and it costs cells.} The measured band is
8.0, 12.1 and 16.1 fine cells for one, two and three levels, against a floor of
five. The margin in the rule is fixed in \emph{base}-cell units, so it is worth
$2^{\text{levels}}$ fine cells. The rule was calibrated by measurement in two
dimensions --- half the margin undershot the floor --- so it is carried over
unchanged here rather than diverging between the two codes, but a margin
expressed in fine cells would recover a factor of two or three in cell count at
three levels.

\missing{the production comparison: same accuracy for how much less}
